\documentclass[11pt]{article}

\usepackage{amsmath}
\usepackage{amssymb}
\usepackage{booktabs}
\usepackage{geometry}
\usepackage{hyperref}
\geometry{margin=1in}

\title{Sample Construction: Mid- and High-Mid-Capitalization Equities\\
with Extended Trading Histories}
\author{}
\date{}

\begin{document}
\maketitle

\section{Data Source}

The universe of candidate securities is drawn from the constituents of the
S\&P 500 index, sourced from the Center for Research in Security Prices
(CRSP) daily stock file. Each security is identified by its CRSP permanent
identifier (\textsc{permno}), which is invariant to ticker or company-name
changes over time. Raw records are restricted to trading dates on or after
January 1, 2000, yielding a panel of $N = 494$ distinct securities.

\section{Longevity Filter}

For each security $i$, let $t_i^{\text{start}}$ and $t_i^{\text{end}}$
denote the first and last observed trading dates in the panel. We define
the observed history length, in years, as
\begin{equation}
  H_i \;=\; \frac{t_i^{\text{end}} - t_i^{\text{start}}}{365.25}.
\end{equation}
A security is retained under the \emph{longevity filter} if
\begin{equation}
  H_i \;\geq\; H_{\min}, \qquad H_{\min} = 20 \text{ years}.
\end{equation}
Applying this filter to the full panel yields $n_{20} = 372$ securities
with at least two decades of continuous trading history.

\section{Market Capitalization Classification}

Let $M_i$ denote the current market capitalization of security $i$, in
billions of USD, obtained from the most recent cross-sectional snapshot of
S\&P 500 constituents. We partition the capitalization axis into four
tiers, following standard practitioner conventions:
\begin{equation}
  \text{Tier}(M_i) =
  \begin{cases}
    \text{Small}     & M_i < 2 \\
    \text{Mid}        & 2 \leq M_i < 10 \\
    \text{High-Mid} \; / \; \text{Low-Large} & 10 \leq M_i < 50 \\
    \text{Large}      & 50 \leq M_i < 200 \\
    \text{Mega}       & M_i \geq 200
  \end{cases}
  \qquad (\text{USD, billions}).
\end{equation}

Within the longevity-filtered set ($n_{20}=372$), the tier distribution is
reported in Table~\ref{tab:tiers}.

\begin{table}[h]
\centering
\caption{Market-capitalization tiers among securities with $H_i \geq 20$ years.}
\label{tab:tiers}
\begin{tabular}{lrr}
\toprule
Tier & Range (USD B) & Count \\
\midrule
Small            & $[0, 2)$    & 0 \\
Mid              & $[2, 10)$   & 14 \\
High-Mid / Low-Large & $[10, 50)$ & 189 \\
Large            & $[50, 200)$ & 130 \\
Mega             & $[200, \infty)$ & 39 \\
\bottomrule
\end{tabular}
\end{table}

The near-total absence of small- and true-mid-cap names in this filtered
set reflects survivorship: firms retaining index membership and continuous
listing for two decades or more have, on average, grown substantially in
capitalization.

\section{Target Universe}

The target universe $\mathcal{U}$ combines the \emph{Mid} and
\emph{High-Mid / Low-Large} tiers, i.e.\ all securities satisfying both the
longevity filter and a capitalization band of \$2B--\$50B:
\begin{equation}
  \mathcal{U} \;=\; \left\{\, i \;:\; H_i \geq 20 \ \text{and} \ 2 \leq M_i < 50 \,\right\}.
\end{equation}
This yields $|\mathcal{U}| = 203$ securities.

\section{Random Partitioning}

To construct $K = 4$ disjoint portfolio sets of equal size $s = 50$, the
following procedure is applied:

\begin{enumerate}
  \item The tickers in $\mathcal{U}$ are randomly permuted using a
        pseudo-random generator with fixed seed ($\text{seed} = 42$), for
        reproducibility.
  \item The first $K \cdot s = 200$ tickers of the permuted sequence are
        retained; the remaining $|\mathcal{U}| - K s = 3$ tickers are
        discarded (AIG, CCI, SBAC).
  \item The retained sequence is partitioned into $K$ contiguous,
        non-overlapping blocks of size $s$:
        \begin{equation}
          \mathcal{U}_k = \{\, u_{(k-1)s+1}, \dots, u_{ks} \,\}, \qquad k = 1,\dots,K,
        \end{equation}
        where $u_1,\dots,u_{200}$ denotes the shuffled ticker sequence.
        By construction, $\mathcal{U}_j \cap \mathcal{U}_k = \emptyset$ for
        $j \neq k$, and $\bigcup_k \mathcal{U}_k$ has no repeated tickers.
\end{enumerate}

It is important to note that randomization is applied only to the
\emph{cross-sectional assignment} of securities to sets. The
\emph{time-series ordering} of each security's own observations is left
untouched: within every output file, rows remain sorted chronologically by
trading date, preserving the temporal dependence structure required for
downstream time-series modeling.

\section{Feature Construction}

For each retained security $i$ and trading day $t$, six predictors are
computed from the raw CRSP fields (price \textsc{prc}, volume
\textsc{vol}, shares outstanding \textsc{shrout}, ask/bid quotes
\textsc{ask}, \textsc{bid}, return \textsc{ret}, and value-weighted market
return \textsc{sprtrn}):

\begin{align}
  \text{RET}_{i,t}          &= \text{(CRSP holding-period return)} \\
  \text{VOL\_CHANGE}_{i,t}  &= \frac{\text{VOL}_{i,t} - \text{VOL}_{i,t-1}}{\text{VOL}_{i,t-1}} \\
  \text{BA\_SPREAD}_{i,t}   &= \frac{\text{ASK}_{i,t} - \text{BID}_{i,t}}{\lvert\text{PRC}_{i,t}\rvert} \\
  \text{ILLIQUIDITY}_{i,t}  &= \frac{\text{RET}_{i,t}}{\text{VOL}_{i,t}\cdot\lvert\text{PRC}_{i,t}\rvert} \\
  \text{sprtrn}_{i,t}       &= \text{(CRSP value-weighted market return)} \\
  \text{TURNOVER}_{i,t}     &= \frac{\text{VOL}_{i,t}}{\text{SHROUT}_{i,t}}
\end{align}

Records with a return code of \texttt{'B'} or \texttt{'C'} (CRSP's markers
for missing or non-standard pricing) are dropped prior to feature
computation, and any resulting infinite values (e.g.\ division by zero
volume) are set to missing.

\section{Output}

The final artifact consists of $K=4$ directories
$\{\texttt{set1},\dots,\texttt{set4}\}$, each containing $s=50$ CSV files
(one per security). Each file contains the columns
$\{\texttt{date}, \text{RET}, \text{VOL\_CHANGE}, \text{BA\_SPREAD},
\text{ILLIQUIDITY}, \text{sprtrn}, \text{TURNOVER}\}$, sorted
chronologically, spanning that security's full available history from
2000 onward.

\end{document}
